% Part: incompleteness
% Chapter: representability-in-q
% Section: basic-representable

\documentclass[../../../include/open-logic-section]{subfiles}

\begin{document}

\olfileid[pt]{inc}{req}{bre}
\olsection{As funções básicas são representáveis em~$\Th{Q}$}

Primeiro, precisamos mostrar que todas as funções básicas são
representáveis em~$\Th{Q}$. Ao final, precisamos mostrar como atribuir
a cada função básica $k$-ária $f(x_0,\dots,x_{k-1})$ !!a{formula}
$!A_f(x_0,\dots,x_{k-1},y)$ que a represente.

Poderemos representar zero, sucessor, adição, multiplicação, a função
característica da igualdade e as projeções. Em cada caso, a fórmula
representante adequada é inteiramente direta; por exemplo, zero é
representada pela fórmula $y = \Obj 0$, sucessor é representada pela
!!{formula} $x_0' = y$ e adição é representada pela !!{formula}
$(x_0 + x_1) = y$. O trabalho consiste em mostrar que $\Th{Q}$ pode
provar as !!{sentence}s pertinentes; por exemplo, dizer que a adição é
representada pela !!{formula} acima envolve mostrar que, para todo par
de números naturais $m$ e $n$, $\Th{Q}$ prova
\begin{align*}
& \eq[\num n + \num m][\num {n+m}] \text{ e}\\
& \lforall[y][(\eq[(\num n + \num m)][y] \lif \eq[y][\num{n+m}])].
\end{align*}

\begin{prop}
\ollabel{prop:rep-zero}
 A função zero $\Zero(x) = 0$ é representada em~$\Th{Q}$ por
$!A_{\Zero}(x,y) \ident \eq[y][\Obj 0]$.
\end{prop}

\begin{prop}
\ollabel{prop:rep-succ}
A função sucessor $\Succ(x) = x+1$ é representada em~$\Th{Q}$ por
$!A_{\Succ}(x,y) \ident \eq[y][x']$.  
\end{prop}

\begin{prop}
\ollabel{prop:rep-proj}
A função de projeção $\Proj{n}{i}(x_0, \dots, x_{n-1}) = x_i$ é
representada em~$\Th{Q}$ por \[!A_{\Proj{n}{i}}(x_0, \dots, x_{n-1}, y) \ident \eq[y][x_i].\]
\end{prop}

\begin{prob}
Prove que $\eq[y][\Obj 0]$, $\eq[y][x']$ e $\eq[y][x_i]$ representam,
respectivamente, $\Zero$, $\Succ$ e $\Proj{n}{i}$.
\end{prob}

\begin{prop}
\ollabel{prop:rep-id}
A função característica de~$=$,
\[
\Char{=}(x_0, x_1) =
\begin{cases}
  1 & \text{se } x_0 =x_1\\
  0 & \text{caso contrário}
\end{cases}
\]
é representada em~$\Th{Q}$ por
\[
  !A_{\Char{=}}(x_0, x_1, y) \ident (\eq[x_0][x_1] \land \eq[y][\num{1}]) \lor (\eq/[x_0][x_1] \land
\eq[y][\num{0}]).
\]
\end{prop}

A prova exige o lema a seguir.

\begin{lem}
\ollabel{lem:q-proves-neq} Dados números naturais $n$ e $m$, se
$n \neq m$, então $\Th{Q} \Proves \eq/[\num n][\num m]$.
\end{lem}

\begin{proof}
Use indução sobre $n$ para mostrar que, para todo $m$, se $n \neq m$,
então $Q \Proves \eq/[\num n][\num m]$.

No caso base, $n = 0$. Se $m$ não é igual a $0$, então $m = k + 1$
para algum número natural $k$. Temos um axioma que diz
$\lforall[x][\eq/[0][x']]$. Por um axioma dos quantificadores,
substituindo $x$ por $\num k$, podemos concluir
$\eq/[0][\num k']$. Mas $\num k'$ é simplesmente $\num m$.

No passo de indução, podemos supor que a afirmação é verdadeira para
$n$ e considerar $n+1$. Seja $m$ um número natural qualquer. Há duas
possibilidades: ou $m = 0$, ou, para algum $k$, temos $m = k+1$. O
primeiro caso é tratado como acima. No segundo, suponha que
$n+1 \neq k+1$. Então $n \neq k$. Pela hipótese de indução para $n$,
temos $\Th{Q} \Proves \eq/[\num n][\num k]$. Temos um axioma que diz
$\lforall[x][\lforall[y][\eq[x'][y'] \lif \eq[x][y]]]$. Usando um
axioma dos quantificadores, temos
$\eq[\num n'][\num k'] \lif \eq[\num n][\num k]$. Usando lógica
proposicional, podemos concluir, em $\Th{Q}$,
$\eq/[\num n][\num k] \lif \eq/[\num n'][\num k']$. Por modus ponens,
podemos concluir $\eq/[\num n'][\num k']$, que é o que queremos, pois
$\num k'$ é $\num m$.
\end{proof}

\begin{explain}
Observe que o lema não diz muito: em essência, afirma que $\Th{Q}$
pode provar que numerais diferentes denotam objetos diferentes. Por
exemplo, $\Th{Q}$ prova $0'' \neq 0'''$. Mas mostrar que isso vale em
geral exige algum cuidado. Observe também que, embora estejamos usando
indução, trata-se de indução \emph{fora} de $\Th{Q}$.
\end{explain}

\begin{proof}[Prova de \olref{prop:rep-id}]
Se $n = m$, então $\num{n}$ e $\num{m}$ são o mesmo termo, e
$\Char{=}(n, m) = 1$. Mas $\Th{Q} \Proves (\eq[\num{n}][\num{m}] \land
\eq[\num{1}][\num{1}])$; portanto, ela prova $!A_=(\num{n}, \num{m},
\num{1})$. Se $n \neq m$, então $\Char=(n, m) = 0$. Por
\olref{lem:q-proves-neq}, $\Th{Q} \Proves \eq/[\num{n}][\num{m}]$
e também $(\eq/[\num{n}][\num{m}] \land \Obj 0 = \Obj 0)$. Logo, $\Th{Q}
\Proves !A_=(\num{n}, \num{m}, \num{0})$.

Para a segunda parte, também temos dois casos. Se $n = m$, precisamos
mostrar que $\Th{Q} \Proves \lforall[y][(!A_=(\num{n}, \num{m}, y)
  \lif \eq[y][\num{1}])]$. Argumentando informalmente, suponha
$!A_=(\num{n}, \num{m}, y)$, isto é,
\[
(\eq[\num{n}][\num{n}] \land \eq[y][\num{1}]) \lor
(\eq/[\num{n}][\num{n}] \land \eq[y][\num{0}])
\]
A disjunção à esquerda implica $\eq[y][\num{1}]$ pela lógica; a da
direita contradiz $\eq[\num{n}][\num{n}]$, que é demonstrável pela
lógica.

Suponha, por outro lado, que $n \neq m$. Então $!A_=(\num{n},
\num{m}, y)$ é
\[
(\eq[\num{n}][\num{m}] \land \eq[y][\num{1}]) \lor
(\eq/[\num{n}][\num{m}] \land \eq[y][\num{0}])
\]
Aqui, a disjunção à esquerda contradiz $\eq/[\num{n}][\num{m}]$, que é
demonstrável em $\Th{Q}$ por \olref{lem:q-proves-neq}; a disjunção à
direita acarreta $\eq[y][\num{0}]$.
\end{proof}

\begin{prop}
\ollabel{prop:rep-add}
A função de adição $\Add(x_0, x_1) = x_0+x_1$ é representada
em~$\Th{Q}$ por
\[
  !A_{\Add}(x_0, x_1, y) \ident \eq[y][(x_0 + x_1)].
\]
\end{prop}

\begin{lem}
\ollabel{lem:q-proves-add}
$\Th{Q} \Proves \eq[(\num{n} + \num{m})][\num{n+m}]$
\end{lem}

\begin{proof}
Provamos isso por indução sobre~$m$. Se $m = 0$, a afirmação é
$\Th{Q} \Proves \eq[(\num{n} + \Obj 0)][\num{n}]$. Isso decorre do
axioma~$!Q_4$. Suponha agora a afirmação para $m$; provemos a afirmação
para $m+1$, isto é, provemos que $\Th{Q} \Proves \eq[(\num{n} +
  \num{m+1})][\num{n+m+1}]$. Observe que $\num{m+1}$ é simplesmente
$\num{m}'$, e $\num{n+m+1}$ é simplesmente $\num{n+m}'$. Pelo axioma
$!Q_5$, $\Th{Q} \Proves
\eq[(\num{n} + \num{m}')][(\num{n}+\num{m})']$. Pela hipótese de
indução, $\Th{Q} \Proves
\eq[(\num{n} + \num{m})][\num{n+m}]$. Logo, $\Th{Q} \Proves
\eq[(\num{n} + \num{m}')][\num{n+m}']$.
\end{proof}

\begin{proof}[Prova de \olref{prop:rep-add}]
A !!{formula}~$!A_\Add(x_0, x_1, y)$ que representa $\Add$ é
$\eq[y][(x_0 + x_1)]$. Primeiro mostramos que, se $\Add(n, m) = k$, então
$\Th{Q} \Proves !A_\Add(\num{n}, \num{m}, \num{k})$, isto é, $\Th{Q}
\Proves \eq[\num{k}][(\num{n} + \num{m})]$. Mas, como $k = n + m$,
$\num{k}$ é simplesmente $\num{n+m}$, e mostramos em
\olref{lem:q-proves-add} que $\Th{Q} \Proves \eq[(\num{n} +
  \num{m})][\num{n+m}]$.

Também precisamos mostrar que, se $\Add(n, m) = k$, então
\[
\Th{Q} \Proves \lforall[y][(!A_\Add(\num{n}, \num{m}, y) \lif
  \eq[y][\num{k}])].
\]
Suponha que tenhamos $\eq[(\num{n} + \num{m})][y]$. Como
\[
\Th{Q} \Proves \eq[(\num{n}+\num{m})][\num{n+m}],
\]
podemos substituir o lado esquerdo por $\num{n+m}$ e obter
$\eq[\num{n+m}][y]$, para~$y$ arbitrário.
\end{proof}

\begin{prop}
\ollabel{prop:rep-mult}
A função de multiplicação $\Mult(x_0, x_1) = x_0 \cdot x_1$ é
representada em~$\Th{Q}$ por
\[
  !A_{\Mult}(x_0, x_1, y) \ident y = (x_0 \times x_1).
\]
\end{prop}

\begin{proof} Exercício. \end{proof}

\begin{lem}
\ollabel{lem:q-proves-mult}
$\Th{Q} \Proves \eq[(\num{n} \times \num{m})][\num{n \cdot m}]$
\end{lem}

\begin{proof} Exercício. \end{proof}

\begin{prob}
Prove \olref[inc][req][bre]{lem:q-proves-mult}.
\end{prob}

\begin{prob}
Use \olref[inc][req][bre]{lem:q-proves-mult} para provar
\olref[inc][req][bre]{prop:rep-mult}.
\end{prob}

\begin{explain}
  Lembre-se de que usamos $\times$ para o símbolo de função da
  linguagem da aritmética e $\cdot$ para a operação usual de
  multiplicação sobre números. Assim, $\cdot$ pode aparecer entre
  expressões de números (como em $m \cdot n$), enquanto $\times$
  aparece apenas entre termos da linguagem da aritmética (como em
  $(\num{m} \times \num{n})$). Para tornar as coisas ainda mais
  confusas, $+$ é usado tanto para a !!{function} quanto para a
  operação de adição. Quando aparece entre termos---por exemplo, em
  $(\num{n} + \num{m})$---, é a !!{function} de $2$ lugares da
  linguagem da aritmética; quando aparece entre números---por exemplo,
  em $n+m$---, é a operação de adição. Isso inclui o caso
  $\num{n+m}$: este é o numeral padrão correspondente ao número~$n+m$.
\end{explain}

\end{document}
